% !TeX root = ../tfg.tex
% !TeX encoding = utf8
%
% ====================================================================
%  APÉNDICE A — PREGUNTAS REALIZADAS A LA EMPRESA Y RESPUESTAS
%
%  RECORDATORIO AL EQUIPO:
%  - Este apéndice es OBLIGATORIO según las instrucciones de la asignatura.
%  - Debéis incluir las preguntas que realizasteis en la entrevista/visita
%    a la empresa y las respuestas que os dieron.
%  - Organizad las preguntas por temática (Tema 2 y Tema 3) o por el
%    orden en que se realizaron.
%  - Podéis incluir el nombre o cargo del entrevistado (si da permiso)
%    y la fecha y medio de la entrevista (presencial, videollamada...).
%  - Si hay información que el entrevistado pidió que no se publicara,
%    indicadlo (p.ej. "Dato confidencial — no se puede reproducir").
%  - Si grabasteis la entrevista (con permiso), podéis transcribir los
%    fragmentos más relevantes.
% ====================================================================

\chapter{Entrevista a CaixaBank: Temas 6, 7 y 8 (Formación, Evaluación y Retribución)}

% ============================================================
\section{Datos de la entrevista}
% ============================================================

\begin{table}[H]
  \centering
  \begin{tabularx}{\textwidth}{lX}
    \toprule
    \textbf{Parámetro}       & \textbf{Detalle} \\
    \midrule
    Empresa                  & CaixaBank \\
    Nombre del entrevistado  & [Anónimo a petición] \\
    Cargo                    & [Director de Recursos Humanos regional] \\
    Departamento             & [Recursos Humanos / Personas y Organización] \\
    Fecha                    & [05/05/2026] \\
    Duración                 & [45 min] \\
    Formato                  & [Videollamada] \\
    \bottomrule
  \end{tabularx}
\end{table}

% ============================================================
\section{Bloque 1: Jesús (Formación y Tendencias Digitales)}
% ============================================================

\subsubsection*{¿Cómo realizan el diagnóstico en los niveles de organización, tareas e individuo para asegurar que la formación mejore la productividad?}

\textbf{Respuesta:}\newline
``A nivel de organización, la estrategia la marca nuestro Plan Estratégico 2025-2027, donde priorizamos, por ejemplo, la IA generativa y la atención al segmento sénior. En el nivel de tareas, nos apoyamos en la Secretaría Técnica de Negocios y en informes de consultoras como Deloitte para definir qué requiere cada puesto (por ejemplo, los perfiles de banca privada requieren analítica compleja). A nivel individual, utilizamos nuestro sistema SAP SuccessFactors para cruzar la evaluación de desempeño del empleado con los requisitos de su puesto actual, saltando alarmas sobre qué formación específica necesita.''

\subsubsection*{Al decidir quién imparte la formación, ¿qué ventajas encuentran en usar formadores internos frente a expertos externos?}

\textbf{Respuesta:}\newline
``Utilizamos un enfoque mixto. El formador interno, como los supervisores o la figura del mentor ('buddy'), es vital para transmitir la cultura organizativa y orientar a nivel práctico sin salir de la sucursal. En cambio, para formación técnica muy específica o que requiere una certificación oficial, como MiFID II o EFPA, nos apoyamos en formadores externos, colaborando con instituciones académicas como la UPF Barcelona School of Management.''

\subsubsection*{¿Qué criterios utilizan para medir el impacto y cómo calculan la relación coste-beneficio?}

\textbf{Respuesta:}\newline
``Medimos el proceso completo: desde la reacción y el aprendizaje mediante cuestionarios y test en nuestra plataforma Evoluciona, hasta la conducta, evaluada por el responsable a los seis o doce meses. Finalmente, el impacto en el negocio lo analizamos usando People Analytics dentro de SAP para ver si esa mejora de capacidades ha incrementado el margen de negocio. Sobre el coste, herramientas como Evoluciona suponen una inversión fuerte inicial, pero al ser e-learning, el coste a la larga y la optimización del tiempo son muchísimo más rentables que la formación presencial tradicional.''

\subsubsection*{¿Cómo están integrando las nuevas tendencias de E-learning, Método 70-20-10 y Universidad Corporativa?}

\textbf{Respuesta:}\newline
``Nuestro pilar es la plataforma Evoluciona, que recientemente ha integrado Microsoft Copilot para crear itinerarios de microaprendizaje ('instant learning') totalmente personalizados. Aplicamos el 70-20-10 de forma muy clara: el 70\% es el día a día en la oficina y la rotación de puestos; el 20\% lo fomentamos a través del mentoring, y el 10\% restante es la formación reglada que damos por Evoluciona, la cual está estructurada funcionalmente a modo de 'escuelas o universidad corporativa'.''

% ============================================================
\section{Bloque 2: Julián (Desarrollo y Planificación de Carrera)}
% ============================================================

\subsubsection*{¿De qué manera el Outdoor training y la Gamificación han mejorado el compromiso de sus empleados?}

\textbf{Respuesta:}\newline
``El outdoor training lo vehiculamos mucho a través del voluntariado corporativo impulsado por la Fundación 'la Caixa', que sirve para cohesionar equipos y vivir los valores del banco fuera de la oficina. En cuanto a la gamificación, la utilizamos directamente en Evoluciona; incluimos retos, puntuaciones y marcadores para hacer atractivos módulos que por norma son más tediosos, como los de prevención de blanqueo de capitales.''

\subsubsection*{¿El plan de carrera de la empresa contempla tanto movimientos verticales como horizontales?}

\textbf{Respuesta:}\newline
``Totalmente. En una red tan extensa, no siempre hay espacio hacia arriba de forma inmediata. Contamos con un sistema de movilidad interna desde el área de People Experience llamado 'Vacantes', que permite traslados entre distintas áreas funcionales. Especialmente tras la integración con Bankia, los movimientos horizontales han sido esenciales.''

\subsubsection*{En las fases de valoración y dirección, ¿qué metodologías utilizan para guiar la trayectoria del empleado?}

\textbf{Respuesta:}\newline
``La valoración parte de una autoevaluación guiada en SAP SuccessFactors, que se complementa con herramientas objetivas de la empresa: evaluación del desempeño de los mandos, centros de evaluación ('assessment centers') para directivos, y planes de sucesión. En la dirección, combinamos servicios de información como el portal CaixaBankCareers con la asesoría profesional de los Business Partners de recursos humanos.''

\subsubsection*{¿En qué situaciones optan por la rotación de puestos, el Coaching, el Mentoring o el Mentoring inverso?}

\textbf{Respuesta:}\newline
``La rotación se fomenta para dar una visión global del negocio a través de 'Vacantes'. El coaching, habitualmente usando el método GROW, lo dirigimos a mandos intermedios y directivos para mejorar el desempeño a corto-medio plazo. El mentoring, en cambio, es una relación a más largo plazo (6 a 12 meses) que usamos con los perfiles de alto potencial, guiados por expertos senior. Finalmente, aunque no tenemos un programa oficial de 'mentoring inverso', la cultura de Evoluciona empuja a que el talento joven nativo digital impulse la transformación tecnológica entre los empleados más veteranos.''

% ============================================================
\section{Bloque 3: Ismael (Gestión de Carrera y Rendimiento)}
% ============================================================

\subsubsection*{¿Cómo incentiva la empresa que el empleado trabaje su Marca Personal, su Networking y su plan de empleabilidad?}

\textbf{Respuesta:}\newline
``Fomentamos las 'carreras múltiples'. El sector está cambiando rápido. Promovemos el networking interno usando redes sociales profesionales, herramientas colaborativas y comunidades de práctica. Queremos que los empleados participen, tengan visibilidad en proyectos transversales y cuiden su empleabilidad fortaleciendo habilidades digitales y de orientación al cliente.''

\subsubsection*{¿Quiénes realizan la evaluación? ¿Se utiliza la evaluación 360º?}

\textbf{Respuesta:}\newline
``La evaluación principal la realiza el responsable o mando directo. Sin embargo, sabemos que en puestos de atención al público o coordinadores, una perspectiva única no basta. Estamos integrando herramientas que recogen feedback más amplio, en la línea de la evaluación 360º, para poder triangular opiniones de clientes internos y del equipo.''

\subsubsection*{¿Qué equilibrio mantienen entre métodos objetivos y subjetivos en la evaluación?}

\textbf{Respuesta:}\newline
``Aplicamos métodos mixtos. Para la parte objetiva pesamos mucho los KPIs de ventas, captación de negocio y NPS de clientes (por ejemplo en un Gestor Comercial suele ser en torno a un 60\%). Y lo combinamos con escalas de valoración cualitativas (el otro 40\%) para evaluar la calidad del servicio, la colaboración y la actitud, con el fin de evitar errores como el efecto halo o evaluar únicamente 'el número' sin importar el 'cómo'.''

\subsubsection*{Ante un bajo rendimiento, ¿qué estrategias de mejora formulan para resolver deficiencias detectadas?}

\textbf{Respuesta:}\newline
``Primero identificamos la raíz: si es falta de capacidad, entra Evoluciona con formación. Si es falta de esfuerzo o desalineación, intervenimos desde el mando con acompañamiento. La meta es no sorprender al final de año, sino ir trazando planes de acción con un seguimiento frecuente.''

% ============================================================
\section{Bloque 4: David (Reforzamiento, Entrevista y Retribución)}
% ============================================================

\subsubsection*{¿Qué importancia le dan al reforzamiento positivo frente al negativo o radical para fomentar conductas deseables?}

\textbf{Respuesta:}\newline
``La cultura de CaixaBank prioriza el refuerzo positivo mediante reconocimientos, méritos y premios a la excelencia ASG (ambientales, sociales y de gobernanza). El sector tiene ya de por sí un refuerzo negativo intrínseco debido a la presión por cumplir objetivos y rankings comerciales. Lógicamente, dejamos las sanciones disciplinarias (refuerzo radical) solo para infracciones graves de fraude o ética operativa.''

\subsubsection*{¿Cómo estructuran la entrevista de evaluación y qué tendencias aplican para medir el rendimiento?}

\textbf{Respuesta:}\newline
``El mánager tiene que hacer un análisis constructivo: revisión de logros, áreas de mejora y alineación de objetivos. En cuanto a las tendencias, hemos cambiado la 'revisión anual' clásica por una evaluación continua, apoyada en reuniones frecuentes ('one-to-one'). Además, especialmente para los líderes, aplicamos evaluación por compromisos ligada a la sostenibilidad y a la inclusión financiera.''

\subsubsection*{¿Cómo es la proporción de la retribución y qué peso tienen los elementos no financieros?}

\textbf{Respuesta:}\newline
``Buscamos retener talento. La retribución directa combina un salario fijo con un bonus comercial variable, que en perfiles directivos se difiere hasta en un 60\% a varios años para asegurar decisiones a largo plazo. Pero nuestro punto fuerte es la retribución indirecta y no financiera: desde nuestro plan de pensiones, seguro médico y aportaciones para ayudas escolares, hasta las políticas de desconexión digital, flexibilidad y programas de voluntariado.''

% ============================================================
\section{Bloque 5: José Ángel (Equidad y Salario Emocional)}
% ============================================================

\subsubsection*{¿Cómo se aseguran de que los salarios base respeten el Convenio Colectivo y cómo gestionan la equidad salarial?}

\textbf{Respuesta:}\newline
``Nos basamos en el IV Convenio de Banca, y de hecho, la mayoría de nuestras retribuciones están muy por encima de las tablas mínimas (un 30-40\% más). La equidad interna la gestionamos a través de 'bandas salariales' en SAP, que clasifican los puestos por su valor. La externa la vigilamos mediante estudios de mercado de consultoras como Mercer, buscando estar en el cuartil superior. Y la equidad individual se plasma con la parte variable o 'bonus' de cada empleado según su desempeño.''

\subsubsection*{¿Cuándo deciden remunerar basándose en el puesto y cuándo en competencias?}

\textbf{Respuesta:}\newline
``Utilizamos un modelo híbrido. El salario base se asigna puramente por el nivel del puesto (Operativo, Gestor, Director de Área...). Pero los complementos y opciones de promoción las basamos en competencias: reconocemos certificaciones externas oficiales (como MiFID II), las habilidades digitales y el liderazgo estratégico.''

\subsubsection*{¿Qué incentivos, Retribución Flexible y Salario Emocional ofrecen?}

\textbf{Respuesta:}\newline
``En incentivos tenemos desde los individuales (el bonus anual o 'stock options') y grupales (objetivos de oficina), hasta participación en beneficios del banco y compra de acciones. Sobre la retribución flexible, el empleado puede contratar desde la nómina seguros médicos, ticket restaurante, ayudas de transporte o guardería, aprovechando ventajas fiscales. Todo ello, combinado con teletrabajo, jornadas intensivas y apoyo psicológico, forma nuestro concepto integral de salario emocional.''

